\chapter{Results}
\label{ch:results}

\pendnote{The $100$-task pilot is running at the time of writing. Every table
and figure in this chapter is laid out with its final shape, its units and its
reading, and every measured value is marked \pend. Nothing in this chapter is
to be read as a result until the marks are replaced.}

\section{Corpus and Measurement Summary}
\label{sec:res-corpus}

Before any comparison, two descriptions of the corpus are needed, because they
decide whether the question RQ1 asks is even well posed. If test costs within a
task were uniform, cost-aware scheduling would have nothing to correct.

\begin{table}[htbp]
\centering
\caption{Admission funnel: how the $100$ pilot tasks were obtained.}
\label{tab:res-funnel}
\small
\begin{tabularx}{\textwidth}{@{}X r r@{}}
\toprule
Stage & Faults remaining & Removed\\
\midrule
ConDefects Python, all faults            & \pend & ---\\
Faulty and reference versions both usable & \pend & \pend\\
At least one test rejects the faulty version & \pend & \pend\\
Reference passes every pooled test        & \pend & \pend\\
Reference finishes within the sandbox clock & \pend & \pend\\
Pool size at or above the minimum         & \pend & \pend\\
\midrule
\textbf{Admitted}                         & \pend & \\
\textbf{Drawn for the pilot}              & $100$ & \\
\bottomrule
\end{tabularx}
\end{table}

\begin{table}[htbp]
\centering
\caption{The two distributions that decide whether cost-aware scheduling has
anything to work with.}
\label{tab:res-arena}
\small
\renewcommand{\arraystretch}{1.15}
\begin{tabularx}{\textwidth}{@{}l r X@{}}
\toprule
Quantity & Value & How to read it\\
\midrule
Tasks / tests / patches & \pend & The size of the pilot universe\\
\midrule
\multicolumn{3}{@{}p{\textwidth}@{}}{\emph{Per-test CPU cost.} If the high
percentiles sit close to the median, costs are near-uniform and the cost term
has nothing to correct. The spread \emph{within} a task is the one that
matters.}\\[2pt]
\quad median              & \pend\,s & \\
\quad 90th percentile     & \pend\,s & \\
\quad 99th percentile     & \pend\,s & \\
\quad median within-task coefficient of variation & \pend & \\
\midrule
\multicolumn{3}{@{}p{\textwidth}@{}}{\emph{Per-test rejection rate.} If a
handful of tests reject everything, any memory finds them quickly and every
configuration converges to the same behaviour.}\\[2pt]
\quad mean                & \pend & \\
\quad median              & \pend & \\
\quad tests that never reject & \pend\,\% & \\
\midrule
\textbf{Median within-task $\mathrm{corr}(\log c,\,K)$} & \textbf{\pend} &
\textbf{The sign decides whether dividing by cost is a refinement or a
correction. If positive, expensive tests are \emph{more} likely to reject, so
ranking on rejection alone is actively drawn towards the expensive end of the
pool.}\\
\bottomrule
\end{tabularx}
\end{table}

\section{RQ1: Effectiveness of Cost-Aware Scheduling}
\label{sec:res-rq1}

\subsection{Head-to-Head}

\begin{table}[htbp]
\centering
\caption[Validation cost per configuration (RQ1)]{RQ1 head-to-head over the
pilot (\code{frtp}~\citep{qi2013frtp}, \code{coverage} =
$\mathrm{RTP_{add}}$~\citep{rothermel2001prioritizing}). The upper block is
absolute cost; the lower block is the paired contrast against \code{sched},
with cluster-bootstrap intervals over tasks ($10{,}000$ resamples) and
Holm-corrected Wilcoxon $p$. \textbf{CPU / exec.} separates ``cheaper'' from
``fewer'': if \code{sched} wins on CPU while paying \emph{more} executions, the
two units of account disagree in sign.}
\label{tab:res-rq1}
\footnotesize
\setlength{\tabcolsep}{5pt}
\begin{tabular}{@{}l r r r r@{}}
\toprule
Configuration & Total CPU (s) & Exec. / round & CPU / round (s) &
\textbf{CPU / exec. (s)}\\
\midrule
\code{natural} (incumbent) & \pend & \pend & \pend & \pend\\
\code{random} (control)    & \pend & \pend & \pend & \pend\\
\code{frtp} & \pend & \pend & \pend & \pend\\
\code{coverage} & \pend & \pend & \pend & \pend\\
\textbf{\code{sched} (ours)} & \pend & \pend & \pend & \pend\\
\midrule
\multicolumn{5}{@{}l@{}}{\emph{Paired contrasts}}\\
\addlinespace
Contrast & $\Delta$ CPU & 95\,\% CI & $p$ / tasks won & $\Delta$ exec.\\
\addlinespace
\code{sched} vs \code{natural}  & \pend & \pend & \pend & \pend\\
\code{sched} vs \code{random}   & \pend & \pend & \pend & \pend\\
\code{sched} vs \code{frtp}     & \pend & \pend & \pend & \pend\\
\code{sched} vs \code{coverage} & \pend & \pend & \pend & \pend\\
\code{frtp} vs \code{natural}   & \pend & \pend & \pend & \pend\\
\bottomrule
\end{tabular}
\end{table}

The contrast against \code{natural} establishes that scheduling is worth doing
at all; the contrast against \code{frtp} tests this work's specific claim,
because \frtp\ already has a memory, so any difference there is attributable to
the rate and the cost term rather than to remembering.

\subsection{Results by Band}
\label{sec:res-bands}

Section~\ref{sec:disc-limits} states three conditions under which the method is
expected to pay off, and Section~\ref{sec:exp-bands} fixed the band boundaries
before the run. Per-band tables for all three axes are produced by
\code{analyse}; the one that carries the argument is
Table~\ref{tab:res-conditions}, which walks from the whole corpus to the
sub-population where all three conditions hold.

\begin{table}[htbp]
\centering
\caption[Whole corpus against the three-condition band]{The two headline views.
The first row is what a deployment gets pointing the tool at everything; the
second is what the mechanism does where the mechanism applies.}
\label{tab:res-conditions}
\footnotesize
\setlength{\tabcolsep}{5pt}
\begin{tabular}{@{}l r r l r l@{}}
\toprule
Population & tasks & \multicolumn{2}{c}{vs \code{natural}} &
\multicolumn{2}{c}{vs \code{frtp}}\\
\cmidrule(lr){3-4}\cmidrule(lr){5-6}
 & & $\Delta$ CPU & 95\,\% CI & $\Delta$ CPU & 95\,\% CI\\
\midrule
Whole corpus & $100$ & \pend & \pend & \pend & \pend\\
\addlinespace
C1 only: $|P| \ge 10$ & \pend & \pend & \pend & \pend & \pend\\
C1 + C2: and costs not flat & \pend & \pend & \pend & \pend & \pend\\
\textbf{C1 + C2 + C3: and not \code{easy}/\code{trivial}} &
  \pend & \pend & \pend & \pend & \pend\\
\bottomrule
\end{tabular}
\end{table}

The expected shape is a gradient, and the gradient is the argument: each step
down removes a sub-population where no effect is mechanically possible. A flat
result across the three steps would mean the conditions of
Section~\ref{sec:disc-limits} are not the operative ones, and the explanation
of \emph{why} the method works is wrong even if the headline number survives.
Two per-band checks are soundness rather than findings: the \code{degenerate}
pool band and the \code{trivial} difficulty band must both show no difference
at all, because in each every configuration runs one identical scan.

\subsection{Learning Curve}

\begin{figure}[htbp]
\centering
\figplaceholder{Mean validation CPU per round, rounds $1$ to $20$, one line per
configuration. Round $1$ is identical across configurations by construction,
because the store is empty.}
\caption{Validation CPU per round. Round~$1$ is the same for every
configuration; what happens afterwards is the store filling in. An incumbent
line that stays flat is the loop never learning anything about its pool.}
\label{fig:res-curve}
\end{figure}

\pendnote{The reading to check: how many rounds the memory configurations need
before they separate from \code{natural}, and whether \code{sched} continues to
separate from \code{frtp} after the rejection information has saturated --- the
cost estimates keep improving after the rejection counts stop changing, so
continued separation is the signature of the cost term doing work.}

\subsection{Soundness Check}

\begin{table}[htbp]
\centering
\caption{Soundness: every Block~A configuration must accept exactly the same
patches on exactly the same rounds (Definition~\ref{def:soundness}). A
difference here is an implementation bug, not a finding.}
\label{tab:res-soundness}
\small
\begin{tabular}{@{}l r r@{}}
\toprule
Configuration & \code{success@B} & First-accept round (mean)\\
\midrule
\code{natural}  & \pend & \pend\\
\code{random}   & \pend & \pend\\
\code{frtp}     & \pend & \pend\\
\code{coverage} & \pend & \pend\\
\code{sched}    & \pend & \pend\\
\bottomrule
\end{tabular}
\end{table}

\section{RQ2: Effectiveness of Selected Hints}
\label{sec:res-rq2}

\begin{table}[htbp]
\centering
\caption[Generation-side results (RQ2)]{Generation-side results. The baseline is \code{no-mem}, which is
independent repeated prompting --- the strategy
\citet{chen2025conversational} found to beat conversational repair;
\code{chat} is a \textsc{ChatRepair}-style transcript~\citep{xia2024chatrepair}.}
\label{tab:res-rq2-main}
\footnotesize
\setlength{\tabcolsep}{5pt}
\begin{tabular}{@{}l r r r r r@{}}
\toprule
Configuration & Dup.\ rate & \code{success@B} & First acc. &
Prompt tok. & Compl.\ tok.\\
\midrule
\code{no-mem} (baseline) & \pend & \pend & \pend & \pend & \pend\\
\code{chat} (transcript) & \pend & \pend & \pend & \pend & \pend\\
\code{exclusion} (prohibition) & \pend & \pend & \pend & \pend & \pend\\
\textbf{\code{hints} (ours)} & \pend & \pend & \pend & \pend & \pend\\
\bottomrule
\end{tabular}
\end{table}

\begin{table}[htbp]
\centering
\caption{Paired contrasts for RQ2, on the duplicate rate.}
\label{tab:res-rq2-contrasts}
\small
\begin{tabular}{@{}l r l r r@{}}
\toprule
Contrast & $\Delta$ duplicate rate & 95\,\% CI & $p$ & tasks won\\
\midrule
\code{hints} vs \code{no-mem}    & \pend & \pend & \pend & \pend\\
\code{hints} vs \code{chat}      & \pend & \pend & \pend & \pend\\
\code{hints} vs \code{exclusion} & \pend & \pend & \pend & \pend\\
\bottomrule
\end{tabular}
\end{table}

Three comparisons run at once and they answer different things. Against
\code{no-mem}, whether showing anything at all helps --- a real bar, since
independent repeated sampling is published as the stronger of the two obvious
strategies. Against \code{chat}, whether \emph{selecting} beats
\emph{accumulating}. Against \code{exclusion}, whether the \emph{form} matters:
it carries the same facts phrased as what not to do, so a difference there
isolates framing from content. The token columns matter because a transcript
grows every round while \code{hints} stays flat by construction --- a
duplicate-rate win that costs twice the tokens is a different result.

\begin{figure}[htbp]
\centering
\figplaceholder{Duplicate rate by round, one line per Block~B configuration,
rounds $1$ to $20$. Round $1$ carries no hints in any configuration.}
\caption{Duplicate rate over the course of a session. The gap between
configurations should widen with round number, because that is when the history
exists to be selected from.}
\label{fig:res-dupes}
\end{figure}

\section{RQ3: Combined Effect}
\label{sec:res-rq3}

\begin{table}[htbp]
\centering
\caption{The $2\times2$ design of Table~\ref{tab:factorial}, measured. The
combined cell is the full system.}
\label{tab:res-rq3}
\footnotesize
\setlength{\tabcolsep}{5pt}
\begin{tabular}{@{}l r r r r@{}}
\toprule
Configuration & Total CPU (s) & Exec. / round & Dup.\ rate &
\code{success@B}\\
\midrule
\code{no-mem} (neither)       & \pend & \pend & \pend & \pend\\
\code{sched} (scheduling only)& \pend & \pend & \pend & \pend\\
\code{hints} (hints only)     & \pend & \pend & \pend & \pend\\
\textbf{\code{sched+hints} (both)} & \pend & \pend & \pend & \pend\\
\bottomrule
\end{tabular}
\end{table}

\begin{table}[htbp]
\centering
\caption{Interaction. ``Additive prediction'' is what the combined cell would
cost if the two mechanisms did not interact; the residual is the interaction
term.}
\label{tab:res-rq3-interaction}
\small
\begin{tabular}{@{}l r r r@{}}
\toprule
Quantity & Additive prediction & Measured & Residual (95\,\% CI)\\
\midrule
Total CPU         & \pend & \pend & \pend\\
Executions / round& \pend & \pend & \pend\\
Duplicate rate    & \pend & \pend & \pend\\
\code{success@B}  & \pend & \pend & \pend\\
\bottomrule
\end{tabular}
\end{table}

The claimed feeding is specific enough to check: hints cut how many redundant
candidates reach validation, which lowers CPU beyond what re-ordering alone
can; scheduling supplies the rejection record hint selection ranks on. A
residual indistinguishable from zero means the two are independent improvements
that happen to ship together --- still useful, but then they are two
contributions rather than one system, and RQ3 is refuted.

\section{RQ4: Ablation and Sensitivity}
\label{sec:res-rq4}

\subsection{Which Term Carries the Scheduling Effect}

\begin{table}[htbp]
\centering
\caption{The ablation ladder of Table~\ref{tab:ladder}, measured. One variable
changes per row. Exploratory: intervals are reported without correction and no
claim rests on a single row.}
\label{tab:res-ladder}
\small
\begin{tabularx}{\textwidth}{@{}l l r l X@{}}
\toprule
Step & What changes & $\Delta$ CPU & 95\,\% CI & Conclusion if it does not clear\\
\midrule
\code{natural} $\to$ \code{cost-only} & cost only        & \pend & \pend & Cost alone is not a usable signal\\
\code{natural} $\to$ \code{kill-only} & rejection only   & \pend & \pend & The memory term is ornament\\
\code{natural} $\to$ \code{frtp}      & published step   & \pend & \pend & \frtp\ does not reproduce here\\
\code{frtp} $\to$ \code{kill-only}    & count $\to$ rate & \pend & \pend & The denominator buys nothing\\
\code{kill-only} $\to$ \code{sched}   & divide by cost   & \pend & \pend & Cost-awareness is ornament\\
\bottomrule
\end{tabularx}
\end{table}

\subsection{How Much to Divide by Cost}

\begin{figure}[htbp]
\centering
\figplaceholder{Total validation CPU against the cost exponent $c$, swept over
$[0, 1.25]$, paired per task with bootstrap bands. $c=0$ is ranking on the
rejection rate alone; $c=1$ is the plain ratio.}
\caption{Cost-exponent sweep. An interior minimum would mean the full cost
division over-corrects, which is what a noisy online denominator predicts; a
minimum at an endpoint would mean one of the two published positions is already
right.}
\label{fig:res-csweep}
\end{figure}

\subsection{The Hint Dose}

\begin{table}[htbp]
\centering
\caption{Hint dose: number of hints crossed with disclosure level. ``Leakage''
is the share of rejected patches whose rejecting tests are disjoint from what
the prompt showed --- the patch satisfied everything it was told and died
elsewhere.}
\label{tab:res-dose}
\small
\begin{tabular}{@{}l l r r r r@{}}
\toprule
$k$ & Disclosure & Duplicate rate & \code{success@B} & Leakage & Prompt tokens\\
\midrule
$1$ & \code{full}            & \pend & \pend & \pend & \pend\\
$2$ & \code{full} (default)  & \pend & \pend & \pend & \pend\\
$3$ & \code{full}            & \pend & \pend & \pend & \pend\\
$2$ & \code{no-required-line}& \pend & \pend & \pend & \pend\\
$2$ & \code{sample-only}     & \pend & \pend & \pend & \pend\\
\bottomrule
\end{tabular}
\end{table}

\begin{figure}[htbp]
\centering
\figplaceholder{Dose-response curve: leakage and \code{success@B} against the
amount disclosed, from \code{sample-only} through \code{no-required-line} to
\code{full}, at $k \in \{1,2,3\}$.}
\caption{The disclosure dose. A curve that is flat across both axes means there
is no dose and the disclosure dial does not exist; a leakage rate that climbs
with disclosure means the model is beginning to satisfy the tests it was shown
rather than repairing the fault.}
\label{fig:res-dose}
\end{figure}

\subsection{Sensitivity to Design Choices}

\begin{table}[htbp]
\centering
\caption{Sensitivity configurations, each paired against \code{sched}. These
belong to threats to validity rather than to the contribution: the question is
whether RQ1's answer survives choices made for other reasons.}
\label{tab:res-sensitivity}
\small
\begin{tabularx}{\textwidth}{@{}l X r l@{}}
\toprule
Configuration & What it changes & $\Delta$ CPU & 95\,\% CI\\
\midrule
\code{sched-cs-median} & Cold start at the median instead of the maximum & \pend & \pend\\
\code{sched-cs-q75}    & Cold start at the $75$th percentile & \pend & \pend\\
\code{sched-bucket}    & Discrete buckets instead of a continuous score & \pend & \pend\\
\code{sched-explore0}  & No exploration ($\varepsilon = 0$) & \pend & \pend\\
\bottomrule
\end{tabularx}
\end{table}

\section{Summary of Findings}
\label{sec:res-summary}

\pendnote{One short paragraph per research question, stating the answer, the
interval, and --- where the pilot does not settle the question --- saying so
plainly rather than reporting a direction that the interval does not support.
RQ1 and RQ3 carry Holm-corrected primary contrasts; the ablation ladder and both
sweeps are exploratory and are to be described as such.}
